\section{Plan testów}
\label{sec:strategia-testow}

Testy podzieliłam na pięć poziomów, zgodnie z ideą piramidy testów. Na dole jest ich najwięcej i są najprostsze, na górze najmniej, ale sprawdzają całą aplikację naraz. Im szerszy test, tym dłużej trwa i tym trudniej znaleźć przyczynę błędu, więc każdą regułę sprawdzałam możliwie nisko.

\begin{table}[htbp]
\centering
\caption{Poziomy testów systemu SZK -- opracowanie własne}
\label{tab:poziomy-testow}
\begin{tabular}{|p{0.15\textwidth}|p{0.26\textwidth}|p{0.31\textwidth}|p{0.09\textwidth}|}
\hline
\textbf{Poziom} & \textbf{Co sprawdza} & \textbf{Narzędzia} & \textbf{Przypadki} \\
\hline
Jednostkowy & Pojedyncze funkcje: obliczenia, walidację, bezpieczeństwo & Vitest & 107 \\
\hline
Integracyjny API & Ścieżki HTTP backendu razem z pośrednikami & Vitest, Supertest, atrapa Supabase & 81 \\
\hline
Integracyjny baza & Polityki Row Level Security i wyzwalacze & Klient HTTP i PostgREST & 16 \\
\hline
Komponentowy & Komponenty i strony aplikacji klienckiej & Vitest, React Testing Library, MSW & 26 \\
\hline
Systemowy & Całe ścieżki użytkownika w przeglądarce & Playwright & 7 \\
\hline
\end{tabular}
\end{table}

Razem wyszły 237 przypadków. 214 z nich działa bez internetu i bez bazy danych, w kilkanaście sekund. Dwa poziomy potrzebują testowej bazy Supabase, więc uruchamiam je osobno.

Przyjęłam trzy zasady. Żaden test nie zależy od sieci, zegara ani losowości -- kurs walutowy i logowanie zastąpiłam atrapami. Wynik oczekiwany liczę sama, ze wzoru albo z treści wymagania; nie wpisuję do testu tego, co akurat zwraca kod, bo wtedy test potwierdzałby błąd zamiast go wykrywać. Nie powtarzam wyżej tego, co sprawdziłam niżej.

Testy API używają atrapy bazy, dzięki czemu mogę wywołać sytuacje trudne do odtworzenia na prawdziwej bazie: błąd zapisu w połowie operacji albo brak profilu użytkownika. Atrapa nie sprawdza jednak reguł zapisanych w samej bazie, więc polityki RLS testuję osobno, wysyłając zapytania prosto do PostgREST. Adres bazy i klucz publiczny widać w kodzie aplikacji klienckiej, więc ktoś może ominąć serwer Express.
